\item \model{Qwen3}

\paragraph*{مقدمه و راهبرد بهینه‌سازی مدل‌های خانواده \model{Qwen3}}
خانواده مدل‌های \model{Qwen3} در بازه مقیاسی ۰.۶ تا ۲۳۵ میلیارد پارامتر در ساختارهای متراکم و \en{MoE} آموزش دیده‌اند \cite{qwen2025qwen3}. این مدل‌ها با معماری ادغام حالت‌های تفکر و بدون تفکر، یکپارچه‌سازی پایداری عددی در پیش‌آموزش و تقطیر بهینه را محقق ساخته‌اند.

\paragraph*{پایدارسازی توجه با \en{QK-Norm} و زیان تعادل بار جهانی}
به منظور تثبیت دینامیک گرادیان و امکان‌پذیری نرخ‌های یادگیری بالا در فازهای آغازین:
\begin{itemize}
    \item بایاس $QKV$ به طور کامل حذف گردیده و عملگر \en{RMSNorm} به صورت مستقل بر بردار کوئری و کلید اعمال می‌شود \cite{dehghani2023scaling}:
    \begin{equation}
        \text{Attn}(Q, K, V) = \text{Softmax}\left( \frac{\text{RMSNorm}(Q) \cdot \text{RMSNorm}(K)^T}{\sqrt{d_k}} \right) V
    \end{equation}
    این رویکرد مانع از واگرایی ضرب داخلی و تولید لاجیت‌های فرین در لایه‌های ابتدایی می‌گردد.
    \item در لایه‌های \en{MoE}، تابع اتلاف تعادل بار در سطح کل دسته توزیع‌شده جهانی (\en{Global-Batch Load Balancing Loss}) \cite{qiu2025demons} اعمال می‌شود:
    \begin{equation}
        \mathcal{L}_{\text{balance}} = \alpha N \sum_{i=1}^N P_i f_i
    \end{equation}
    که در آن $P_i$ میانگین احتمال تخصیص متخصص $i$ و $f_i$ نسبت کل توکن‌های هدایت‌شده به آن در سراسر پردازنده‌ها است.
\end{itemize}

\paragraph*{پایپ‌لاین چهارمرحله‌ای پس‌آموزش در مدل‌های پرچم‌دار}
\begin{itemize}
    \item \textbf{مرحله ۱: آغاز سرد استدلال (\en{Long-CoT Cold Start}):} آموزش با داده‌های استدلالی تصفیه‌شده جهت پی‌ریزی الگوهای تفکر چندمرحله‌ای.
    \item \textbf{مرحله ۲: یادگیری تقویتی استدلال (\en{Reasoning RL}):} بهینه‌سازی با \en{GRPO} روی ۳,۹۹۵ جفت پرسش-ارزیاب با دسته‌های بزرگ و حفظ پویای آنتروپی.
    \item \textbf{مرحله ۳: همجوشی حالت‌های تفکر (\en{Thinking Mode Fusion}):} آموزش پیوسته \en{SFT} با تلفیق پرومپت‌های دارای \verb|/think| و \verb|/no_think| جهت سازگارسازی مدل با تصمیم‌گیری در شرایط توقف اضطراری تفکر.
    \item \textbf{مرحله ۴: یادگیری تقویتی عمومی (\en{General RL}):} ارتقای کیفیت در حوزه‌های مکالمه، پیروی از دستورات و استفاده از ابزار با ترکیبی از پاداش‌های ارزیاب، مبتنی بر مدل و مدل‌های مولد پاداش.
\end{itemize}

\paragraph*{تقطیر خط‌مشی قوی به ضعیف (\en{Strong-to-Weak Distillation})}
برای مدل‌های سبک‌تر (۰.۶B تا ۱۴B و ۳۰B-A3B)، کل چرخه چهارمرحله‌ای با تقطیر لاجیت از مدل‌های پرچم‌دار جایگزین می‌شود:
\begin{equation}
    \mathcal{L}_{\text{distill}}(\theta) = \mathbb{E}_{x \sim \mathcal{D}, y \sim \pi_\theta} \left[ D_{\text{KL}}(\pi_\theta(\cdot \mid x, y_{<t}) \parallel \pi_{\text{teacher}}(\cdot \mid x, y_{<t})) \right]
\end{equation}
این رویکرد با مصرف تنها یک‌دهم منابع یادگیری تقویتی، توانایی استدلال و قابلیت تعویض حالت را به مدل‌های کوچک منتقل می‌سازد.

\begin{figure}[htbp]
\centering
\begin{tikzpicture}[
    node distance=1.1cm and 1.8cm,
    stage/.style={draw=blue!70!black, fill=blue!10, rounded corners, align=center, font=\small, text width=5cm, line width=1pt},
    arrow/.style={-{Stealth[scale=1.1]}, line width=1pt, draw=gray!80!black}
]
    \node (s1) [stage] {\textbf{مرحله ۱: آغاز سرد استدلال}\\آموزش با زنجیره تفکر عمیق (\en{Long-CoT})};
    \node (s2) [stage, below=of s1] {\textbf{مرحله ۲: یادگیری تقویتی استدلال}\\بهینه‌سازی \en{GRPO} بر مسائل تحلیلی};
    \node (s3) [stage, below=of s2] {\textbf{مرحله ۳: همجوشی حالت تفکر}\\تلفیق حالات با پرومپت \en{/think} و \en{/no\_think}};
    \node (s4) [stage, below=of s3] {\textbf{مرحله ۴: یادگیری تقویتی عمومی}\\بهینه‌سازی جامع با پاداش‌های سه‌گانه};
    
    \draw [arrow] (s1) -- (s2);
    \draw [arrow] (s2) -- (s3);
    \draw [arrow] (s3) -- (s4);
\end{tikzpicture}
\caption{پایپ‌لاین چهارمرحله‌ای پس‌آموزش و همجوشی حالت‌های تفکر در \model{Qwen3}}
\label{fig:qwen3_post_train}
\end{figure}